\chapter{Firmas enteras y valores numéricos de referencia}
\label{app:tablas}

\section{Restricción histórica calibrada de diez filas}

La tabla siguiente define la restricción calibrada \(\mathcal V_{10}\) utilizada
en \(E_0\)--\(E_3\). Tanto las firmas como el vector de energías equivalentes
\(E_{\rm ref}=m_{\rm ref}c^2\) son datos congelados de esta evaluación y se
reproducen literalmente en
\path{datos/tabla_etapas_caracter_masa.csv}. No se atribuye ese vector a una
edición posterior del Particle Data Group; las comparaciones con una edición
externa deben citarla por separado \cite{PDG2024}.
Esta tabla es una reconstrucción calibrada condicionada a esos objetivos; no
constituye una predicción prospectiva en ausencia de un objetivo
\(E_X^{\rm ref}\).

\begin{longtable}{@{}lrrrrrr@{}}
\caption[Firmas y energías equivalentes de referencia]{Firmas de acción de
\(\mathcal V_{10}\) y energías equivalentes de
referencia \(E_X^{\rm ref}=m_X^{\rm ref}c^2\), en MeV.}
\label{tab:firmas-diez}\\
\toprule
Especie & \(n_A\)&\(s_C\)&\(k\)&\(\nu_{120}\)&\(\nu_{270}\)&\(E_{\rm ref}\)\\
\midrule
\endfirsthead
\toprule
Especie & \(n_A\)&\(s_C\)&\(k\)&\(\nu_{120}\)&\(\nu_{270}\)&\(E_{\rm ref}\)\\
\midrule
\endhead
\(\mu\)&42&-3&8&0&2&105.6583755\\
\(\pi^0\)&44&-4&-25&1&-2&134.9768\\
\(\pi^+\)&44&1&-2&2&-1&139.57039\\
\(K^+\)&54&-1&17&-2&2&493.677\\
\(K^0\)&54&0&32&3&-2&497.611\\
\(p\)&59&0&9&1&0&938.27208816\\
\(n\)&59&1&-34&0&1&939.56542052\\
\(\tau\)&64&0&25&-1&6&1776.86\\
\(W\)&94&-1&0&-2&4&80379.0\\
\(Z\)&95&-1&-11&0&-4&91187.6\\
\bottomrule
\end{longtable}

\section{Evaluación sucesiva de la acción}

La raíz cuadrática media de los errores relativos se abrevia RMS, por
\emph{root mean square}.

\begin{table}[htbp]
\centering
\caption{RMS relativo en las etapas sucesivas de la acción.}
\label{tab:rms-etapas}
\begin{tabular}{@{}lrr@{}}
\toprule
Etapa & escala electrónica HMT & escala externa de calibración\\
\midrule
\(E_0\)&\(2.2242372039\times10^{-3}\)&\(2.2242321491\times10^{-3}\)\\
\(E_1\)&\(2.7620461148\times10^{-5}\)&\(2.7610767376\times10^{-5}\)\\
\(E_2\)&\(4.7772213690\times10^{-6}\)&\(4.7655247180\times10^{-6}\)\\
\(E_3\)&\(4.7953121368\times10^{-7}\)&\(4.7557629976\times10^{-7}\)\\
\(E_4\)&\(2.2292793330\times10^{-8}\)&\(6.5144600068\times10^{-8}\)\\
\bottomrule
\end{tabular}
\end{table}

Las dos escalas electrónicas empleadas, expresadas como energías equivalentes,
son
\[
 E_e^{\rm HMT}=m_e^{\rm HMT}c^2
 =0.510998922488066\,\mathrm{MeV},
 \qquad
 E_e^{\rm cal}=m_e^{\rm cal}c^2
 =0.510998950000000\,\mathrm{MeV}.
\]
El estadio \(E_4\) es la cuantización
inversa del residuo posterior a \(E_3\), con
\[
 \nu_{180}=(2,7,-7,-1,-4,5,-4,-5,6,9).
\]

\section{Sensibilidad respecto de un segundo vector de comparación}

El fichero reproducible contiene un segundo vector numérico, congelado como
ensayo de sensibilidad y no atribuido aquí a una edición metrológica concreta.
Este vector cambia de forma apreciable los valores de \(\tau,W,Z\). En el
conjunto completo de diez filas históricas, los RMS relativos son
\[
 \operatorname{RMS}_{\rm post}(E_3)=6.6028250804\times10^{-5},
\]
\[
 \operatorname{RMS}_{\rm post}(E_4)=6.6212121422\times10^{-5}.
\]
En las siete especies que excluyen \(\tau,W,Z\), son
\[
 6.4172392004\times10^{-7},
 \qquad
 3.6764316632\times10^{-7},
\]
respectivamente. La diferencia cuantifica la sensibilidad de la evaluación a
la elección de referencia y exige identificarla en toda comparación.

\section{Coordenadas angulares y observables derivados}

\begin{table}[htbp]
\centering
\caption{Constantes y observables angulares utilizados en los cálculos.}
\label{tab:resumen-angular}
\begin{tabular}{@{}ll@{}}
\toprule
Objeto&Valor\\
\midrule
\(A\)&\(7.297352569283801^\circ\)\\
\(\Cstar\)&\(2.123738338968646^\circ\)\\
\(h=\Cstar/6\)&\(0.353956389828107621\ldots{}^\circ\)\\
\(\Delta\)&\(0.00637108571721236^\circ\)\\
\(\Gamma_{6\to8}\)&\(0.274007672694459275\ldots\)\\
\(\theta_{12}\)&\(13.0033135895090^\circ\)\\
\(\theta_{23}\)&\(2.39805615733160^\circ\)\\
\(\theta_{13}\)&\(0.213977382243507^\circ\)\\
\(\delta_{\rm CP}\)&\(65.6761731235542^\circ\)\\
\(J\)&\(3.1189723326632974\times10^{-5}\)\\
\bottomrule
\end{tabular}
\end{table}

Los apéndices digitales conservan la expansión decimal completa empleada en
estas evaluaciones; las tablas presentan las cifras significativas necesarias
para la lectura del resultado.
